\section{Last Year's Midterm}

\begin{enumerate}[label=\textbf{\arabic*.}]
    \item
    \begin{enumerate}[label=(\alph*)]
        \item True. Since \(x^TA^TAx=\|Ax\|^2\), if \(A^TA=0\), every
        column of \(A\) has zero norm.
        \item True. \(A^TA\) is symmetric. A symmetric upper triangular
        matrix must be diagonal.
        \item False. For example, \(A=2I\) gives \(A^TA=4I\), which is
        diagonal, but \(A\) is not orthogonal.
        \item False. For example,
        \[
            A^TA=
            \begin{bmatrix}
                1&-1\\
                -1&1
            \end{bmatrix}
        \]
        is possible and has all off-diagonal entries negative.
    \end{enumerate}

    \item \(S=\operatorname{range}(U)\). Since \(U^TU=I\), \(UU^T\) is
    the orthogonal projector onto \(\operatorname{range}(U)\).

    \item \(S\cup T\) is a subspace only if \(S\cap T=S\). A union of two
    subspaces is a subspace only when one subspace is contained in the
    other; since \(\dim(S)<\dim(T)\), this means \(S\subseteq T\).

    \item True. If \(ABx=0\), then \(Bx\in\operatorname{null}(A)\), so
    \(Bx=0\). Since \(\operatorname{null}(B)=\{0\}\), \(x=0\).

    \item False. Let
    \[
        A=\begin{bmatrix}1&0\end{bmatrix},
        \qquad
        B=\begin{bmatrix}0&0\\0&1\end{bmatrix}.
    \]
    Then \(\operatorname{null}(A)\cap\operatorname{null}(B)=\{0\}\), but
    \(AB=0\).

    \item False. Additivity alone does not imply homogeneity over all
    real scalars for an arbitrary function.% Without an additional
    % regularity assumption, additive non-linear functions can exist.

    \item True. Take \(U\) to be rotation by angle \(2\pi/n\). Then
    \(U^n=I\), and no smaller positive power is the identity.

    \item False. The recursive construction is missing the factor
    \(1/\sqrt{2}\) at each step. Without it, the block matrix has columns
    with norm \(\sqrt{2}\), not \(1\).

    \item
    \begin{enumerate}[label=(\alph*)]
        \item False. For example, \(u=v=e_1\) gives
        \(A=e_1e_1^T\), which is upper triangular.
        \item True. The sum of the diagonal entries is
        \(
            \operatorname{trace}(uv^T)=v^Tu=u^Tv.
        \)
        \item \(Q\) and \(u\) are parallel but may point in opposite
        directions. Each nonzero column of \(A\) is a scalar multiple of
        \(u\), and the sign depends on the scalar from \(v\).
    \end{enumerate}

    \item True. \(n-1\) independent vectors in \(\mathbb{R}^n\) span an
    \((n-1)\)-dimensional subspace, which is a hyperplane.

    \item False. The set is the points closer to \(b\) than to the
    origin:
    \[
        \|x-b\|_2\leq \|x\|_2
        \Longleftrightarrow
        b^Tx\geq \frac{1}{2}\|b\|_2^2
    \]
    This is a halfspace. If \(b\neq 0\), the origin is not in the set,
    since \(b^T0=0<\frac{1}{2}\|b\|_2^2\). Therefore the set is not a
    subspace.

    \item True. If a solution exists, then there are infinitely many,
    since \(\dim\operatorname{null}(A)\geq n-m>0\). Otherwise there are
    no solutions.

    \item
    \begin{enumerate}[label=(\alph*)]
        \item \(A\Leftrightarrow B\). The rank of \(\begin{bmatrix}A&B\end{bmatrix}\)
        equals the rank of \(B\) exactly when the columns of \(A\) lie in
        \(\operatorname{range}(B)\).
        \item \(A\not\Leftrightarrow B\). A vector outside
        \(\operatorname{range}(A)\) need not be orthogonal to the range,
        and \(y=0\) is in \(\operatorname{range}(A)^\perp\) but also in
        \(\operatorname{range}(A)\).
        \item \(A\Rightarrow B\). Infinitely many solutions imply a
        nonzero nullspace direction. The converse can fail if \(y\) is
        not in \(\operatorname{range}(A)\).
    \end{enumerate}

    \item Interpret \(T(h)\) as the dimension of the span of the range
    of \(h\). Since \(f\) and \(g\) are linear maps, this is the rank of
    the corresponding matrix.
    \begin{enumerate}[label=(\alph*)]
        \item False. For example, take
        \[
            F=
            \begin{bmatrix}
                0&0\\
                0&1
            \end{bmatrix},
            \qquad
            G=
            \begin{bmatrix}
                0&0\\
                1&0
            \end{bmatrix}
        \]
        Then
        \[
            FG=
            \begin{bmatrix}
                0&0\\
                1&0
            \end{bmatrix},
            \qquad
            GF=
            \begin{bmatrix}
                0&0\\
                0&0
            \end{bmatrix}
        \]
        so \(\operatorname{rank}(FG)=1\), but
        \(\operatorname{rank}(GF)=0\).
        \item True. \(\operatorname{rank}(F^2)\leq\operatorname{rank}(F)\),
        and \(\operatorname{rank}(F+F)=\operatorname{rank}(2F)=\operatorname{rank}(F)\).
        \item False. \(\operatorname{rank}(FG)\leq\operatorname{rank}(G)\).
        \item True. Choose \(g\) to be the identity map. Then
        \(f\circ g=g\circ f=f\).
        \item False. If \(f=0\), then \(T(f\circ g)=0\), so it cannot
        equal \(T(g)>0\).
        \item False. The same
        \(G=\begin{bmatrix}0&0\\1&0\end{bmatrix}\) is nonzero but
        nilpotent:
        \[
            G^2=0
        \]
        \item True. If \(T(f)\geq T(g)\) for every \(g\), then \(f\) has
        full rank and is invertible, so \(T(f\circ f)=T(f)\).
        \item True. When \(m>n\), such an \(f\) has full column rank, so
        it has a left inverse \(h\) with \(h(f(x))=x\).
        \item False. A full-column-rank map \(f:\mathbb{R}^n\to
        \mathbb{R}^m\) with \(m>n\) is not onto \(\mathbb{R}^m\), so it
        cannot have a right inverse onto all of \(\mathbb{R}^m\).
        \item False. If \(T(f)\leq T(g)\) for every \(g\), then \(f=0\).
        Hence \(\|f(x)\|=0\), never greater than \(\|x\|\).
    \end{enumerate}
\end{enumerate}
